\section{Escopo da Prova 1 e o que cada seção cobra}

A lista da professora (tabela de questões) vai de 4.9 até 6.3 antes de mudar de formato (7.1 em diante é ``aquecimento / lista completa'', material da P2). O cronograma das aulas 02--05 detalha as primeiras seções. Cruzando os dois:

\begin{center}
\small
\renewcommand{\arraystretch}{1.25}
\begin{tabular}{@{}p{1.0cm}>{\raggedright\arraybackslash}p{2.9cm}>{\raggedright\arraybackslash}p{3.4cm}>{\raggedright\arraybackslash}p{8.3cm}@{}}
\toprule
\textbf{Seção} & \textbf{Tema} & \textbf{Questões recomendadas} & \textbf{Tipo de questão que isso gera na prova}\\
\midrule
4.9 & Primitivas & 1--26; 29--54; 55; 56; 65--70 & primitiva mais geral; achar $f$ a partir de $f'$ ou $f''$ com condições; reta tangente dada; cinemática $a\to v\to s$\\
5.1 & Áreas e distâncias & 2, 3, 11, 20, 24 (aula 04); 16--23 & estimar área com $L_n$, $R_n$, $M_n$ e dizer se é sub/superestimativa; escrever a área como limite; reconhecer a região dada por um limite\\
5.2 & Integral definida & 1, 3, 19--22, 23, 25--34, 41--46, 49, 50 & calcular soma de Riemann com $n$ dado \emph{e dizer o que ela representa}; limite $\leftrightarrow$ integral; calcular pela definição (Teorema 4); calcular por áreas\\
5.3 & TFC & 9, 10, 21--24, 25--54, 85, 86 & derivar $\int_a^{x} f(t)\dt$ (com regra da cadeia); calcular pelo TFC2 e interpretar como área; reconhecer um limite como soma de Riemann em $[0,1]$\\
5.4 & Integral indefinida e Variação Total & 1--24; 27--54; 57; 59--62; 69--74 & verificar fórmula por derivação; integral indefinida geral; definidas; interpretar $\int$ taxa; deslocamento $\times$ distância; massa, vazão\\
5.5 & Substituição & 1--8; 9--54; 59--80; 86--88 & substituição $u$ em indefinidas e definidas (trocando limites); simetria par/ímpar\\
6.1 & Área entre curvas & 1--6; 11--25 & esboçar, achar interseções, integrar em $x$ ou em $y$\\
6.2 & Volumes & 1--40 & discos e anéis; eixos deslocados ($y=k$, $x=k$)\\
6.3 & Cascas cilíndricas & 2--6; 9; 10; 25--30 & cascas; escolher entre cascas e anéis\\
\bottomrule
\end{tabular}
\end{center}

\begin{dica}[{O que \emph{não} cai na P1 --- e o que isso muda}]
Integração por partes, frações parciais e substituição trigonométrica são do capítulo 7 (P2). Portanto, na P1 a ``classificação da técnica'' se reduz a quatro opções:
\begin{enumerate}[leftmargin=1.6em, itemsep=1pt]
\item \textbf{Tabela direta} após preparação algébrica (expandir, dividir termo a termo, identidade trigonométrica) --- seções 4.9, 5.3, 5.4;
\item \textbf{Substituição $u$} quando há função composta acompanhada da derivada ``de dentro'' --- seção 5.5;
\item \textbf{Interpretação geométrica} (retângulo, triângulo, trapézio, semicírculo, valor absoluto) --- seção 5.2;
\item \textbf{Definição} como limite de somas de Riemann, quando o enunciado manda --- seções 5.1, 5.2.
\end{enumerate}
Se você se pegar tentando ``por partes'' na P1, pare: a questão está pedindo uma das quatro acima.
\end{dica}
